\subsection{MIND-small proxy construction}
\label{app:mind_proxy}
MIND provides click histories and impression slates rather than independent stated profiles. We therefore construct a public proxy with the blocks in Table~\ref{tab:mind_proxy_fields}. Here P summarizes the history as a profile, whereas C summarizes the candidate's similarities to individual history articles. Their computations differ, but both originate from prior clicks.

\begin{table}[H]
\centering\small
\caption{MIND proxy input blocks. Counts are scalar fields per candidate, not embedding dimensions. These public-data definitions differ from the private 25-field schema.}
\label{tab:mind_proxy_fields}
\setlength{\tabcolsep}{5pt}
\renewcommand{\arraystretch}{1.12}
\begin{tabularx}{\textwidth}{@{}p{0.10\textwidth}rY>{\raggedright\arraybackslash}p{0.18\textwidth}@{}}
\toprule
Block & Fields & Construction & Source \\
\midrule
P proxy & 3 & Cosine alignment with the normalized mean-history vector; fractions of retained history articles matching the candidate's category and subcategory. & Prior click history \\
E & 3 & Logged-slate membership, reciprocal recorded position, and position percentile; zero for added candidates. & Recorded impression slate \\
C & 4 & Mean, maximum, top-three mean, and attention-weighted mean of candidate-to-history cosine similarities. & Prior click history \\
Controls & 8 & Article, category, and subcategory click-through-rate estimates; log-transformed article impression count, history size, slate size, title length, and abstract length. & Training logs and metadata \\
Products & 6 & $P_0E_m$, $C_{\mathrm{mean}}E_m$, $C_{\mathrm{attn}}E_m$, $C_{\max}E_m$, $P_0C_{\mathrm{mean}}$, and $P_0C_{\mathrm{attn}}$, where $P_0$ is profile cosine alignment and $E_m$ is slate membership. & Derived P/E/C fields \\
\bottomrule
\end{tabularx}
\end{table}

\noindent\textbf{Model inputs.} Controls+E uses $8+3=11$ fields; P+C uses $8+3+4=15$; the full PEC proxy uses $8+3+3+4+6=24$. Thus the P+C--full comparison changes both E and the six products, not E alone. The artifact key \texttt{exposure\_only} denotes Controls+E, not the private trace-only diagnostic in Appendix~\ref{app:e_only_diagnostic}.

\noindent\textbf{Feature calculation.} Article vectors encode title and abstract and are L2-normalized. Histories retain up to 50 preceding clicked articles with available vectors. P alignment normalizes the mean history vector before taking its cosine with the candidate; $C_{\mathrm{mean}}$ instead averages individual cosines, so the two values need not coincide. C attention uses softmax similarity weights with temperature 0.18, without the private recency weighting. Empty histories yield zero P/C values; top-three uses all available similarities when fewer than three exist. Training click-through-rate estimates use $(\mathrm{clicks}+1)/(\mathrm{impressions}+2)$.

\noindent\textbf{Candidate construction.} Each impression defines one query. Logged candidates retain their click labels and recorded slate positions. Up to ten additional articles are sampled without replacement from eligible items in the 5,000 most frequently impressed training articles, excluding the logged slate and retained history. Sampling uses seed 42 with a query-identifier-derived random generator. Added items receive label zero as weak negatives, not verified dislike or non-exposure. Both P and C remain history-derived, so this audit assesses the public proxy rather than reproducing independent stated-profile provenance.
